\chapter{聚类分析}
\chaptergoal{理解聚类是基于所定义相似性和结构规则的无监督任务；掌握 K-means、层次聚类与 DBSCAN 的基本机制；比较参数、尺度、初始化和噪声对结果的影响；能够结合 Gaia 或其他外部信息检查聚类候选，而不把聚类编号直接当作真实类别。}

聚类尝试在没有训练标签的条件下，从数据中发现结构。人们常把同一片天空中的星群或相近物理测量当作潜在群体，但算法输出只回答“在指定特征和距离规则下，这些对象呈现怎样的结构”。要把结构进一步解释为真实星团、材料族或物理状态，还需要独立证据。
\section{无标签数据中的结构问题}

\textbf{簇（cluster）}是聚类算法按其规则形成的一组样本。它与监督分类中的\textbf{类别（class）}不同：类别通常由已有参考标签定义，簇则由所选特征、距离和算法决定。因此同一数据在不同特征或规则下可以有不同而各自有用途的分组。

\textbf{相似性}必须落实到一个可计算标准。例如只使用花瓣长度时，两朵花长度接近就可能被分到一起；再加入萼片宽度，邻近关系就可能改变。本章先讨论算法定义的结构，随后才讨论这种结构是否具有专业含义。

假设一组二维点形成三个相隔较远的团块，把它们分为三个群体似乎十分自然；若数据形成两条弯曲的月牙，按到固定中心的距离划分就可能不适合。聚类方法的设计包含“什么叫一个群体”的假设，K-means 偏向紧凑、相对球形的群体，密度方法则可寻找不规则形状的连通高密度区域。

训练时不使用真实类别标签，不意味着我们不能在实验后借助标签检查结果。但必须把“事后对照”和“使用标签确定聚类参数”区分开；若用标签选择聚类方案，该过程就引入了监督信息，不能再把整体流程描述为完全未使用标签。

\subsection{没有标签，不等于没有目标函数}
同一组点可以按空间距离、光谱形状或运动特征分组，不同特征与距离会产生不同群体。算法不会自动知道我们关心的是星团成员、仪器状态还是材料族。聚类之前必须说明对象、特征、相似关系和希望支持的分析。

聚类的输出还可能有不同形式：每条记录分到一个簇；允许某些记录标为噪声；或对多个群体给出软分配。簇编号只是计算标识，交换0和1不改变分组。评价稳定性时应比较成员关系或采用对编号置换不敏感的指标，而不是逐个检查簇号是否相同。

聚类后使用参考类别比较是有效分析，但若反复根据这些类别调整参数，整个流程已经含有监督选择信息。可以这样设计研究，但不能仍声称参数完全在无标签条件下确定。本书先冻结生成实验的设置，再用已知生成标签解释结果。

\section{K均值的中心与迭代}

\textbf{K均值（K-means）}中的$K$是预先指定的簇数，“均值”是各簇样本坐标的平均值，也称\textbf{质心}。质心不一定恰好是一条实际记录。算法先给出一组初始中心，再反复执行“按最近中心分组”和“按分组更新均值”。

用四个人工一维点$0,2,8,10$作完整演示，设$K=2$，初始中心是0和2。第一轮按距离分配：0进入第一簇，2、8、10进入第二簇。新的中心是0和$(2+8+10)/3=20/3$。第二轮分配时，0和2更靠近0，8和10更靠近$20/3$；更新中心变成1和9。再分配一次，分组不变，达到这次运行的稳定结果。最终簇内平方和为
\[
(0-1)^2+(2-1)^2+(8-9)^2+(10-9)^2=4.
\]
若只设一个簇，中心为5，平方和为68；若设四个簇且每点自成一簇，平方和为0。可见只看这个训练目标会偏好更多簇，不会自动告诉我们“正确簇数”。这里的中心更新是按分配求均值，不是第3章的逐参数梯度下降。

K-means 预先指定簇数 $K$，交替执行两步：将每个样本分配给最近的中心；重新计算每个簇的中心。其目标常写为
\begin{equation}
J=\sum_{k=1}^K\sum_{i\in C_k}\|\mathbf x_i-\boldsymbol\mu_k\|^2,
\end{equation}
$C_k$ 是第 $k$ 个簇，$\boldsymbol\mu_k$ 是中心。算法逐步改善这个目标，但不同初始化可能得到不同的局部结果。数据尺度、异常点和簇大小也会影响最终划分。

选择 $K$ 不能只看训练目标：簇数增大时，簇内平方距离通常不会上升。可观察肘部趋势、轮廓系数、结果稳定性和外部科学约束，但这些指标各有局限，没有一个指标能自动告诉我们真实天体群体的数量。

\subsection{从目标函数理解分配与更新}
设要分成 $K$ 个簇，中心为 $\boldsymbol\mu_1,\ldots,\boldsymbol\mu_K$，样本归属为 $c_i$。K均值最小化簇内平方距离
\[
J=\sum_i\|\mathbf x_i-\boldsymbol\mu_{c_i}\|^2.
\]
若中心固定，每条记录分给最近中心使该项最小；若归属固定，对每个簇求和，最优中心是该簇均值，与第8章平方误差下叶值为均值的证明相同。Lloyd迭代交替执行两步，每一步在当前另一部分固定时不增大目标。

不增大目标不等于找到全局最优。分配关系是离散的，不同初始中心可以落入不同稳定解；一旦所有记录保持归属、中心也不变，算法可以停止在局部固定点。出现空簇时还需要规定重新初始化等处理，不能对空集合直接求均值。

\worked{一次迭代可以稳定，却仍有更好的分组}
取一维人工点 $0,1,4,5,9,10$，$K=2$，初始中心0和4。按最近中心分配，前两个点归第一簇，其余归第二簇；更新后中心为0.5和7。此时中点3.75，下一轮分配不变，簇内平方和为
\[
(0-0.5)^2+(1-0.5)^2+(4-7)^2+(5-7)^2+(9-7)^2+(10-7)^2=26.5.
\]
换一组候选中心2.5和9.5，分配为 $\{0,1,4,5\}$ 与 $\{9,10\}$，平方和为 $17+0.5=17.5$，更低。前一个算法确实收敛了，但收敛不是全局最优的证明。配套程序保存每轮中心和目标，学生可以另用不同初值验证这一差异。

\subsection{初始化、重复运行与K的选择}
K-means++通过相对分散的候选初始化改善起点，仍不保证所有数据上找到全局最优。\texttt{n\_init}控制尝试多个初始化的次数，选择目标较小者；不同软件版本的默认可能改变，教学代码应显式设置，例如 \texttt{n\_init=10}，不依赖隐含配置。

增加 $K$通常能降低可达到的训练平方距离，极端 $K=n$时每条记录单独成簇就为0。因而不能仅靠最小训练目标选择 $K$。肘部图寻找边际改善减弱的位置，是启发式而不是严格证据；轮廓系数、重采样稳定性和专业可解释性提供不同信息，可能给出不同偏好。

K均值倾向于形成在欧氏距离下围绕中心较紧凑的区域。非凸弯曲形状、不平衡密度或尺度不同的群体不一定符合这个目标。知道目标函数，就能预判为什么某些形状会被算法切开，而不必把它解释为软件错误。

\section{层次聚类的基本思想}

\textbf{凝聚式层次聚类}开始时让每个点自成一簇，每一步合并当前最接近的两簇，记录合并距离，从而得到多层分组关系。\textbf{树状图}把这些合并及其距离画成树，图中的高度表示相应连接准则下的合并尺度，而不是一个分类树的预测阈值。

例如一维点0、1、4，先合并最近的0与1，得到簇A=$\{0,1\}$和B=$\{4\}$。\textbf{单链接}取两簇间最近点对距离，得3；\textbf{完全链接}取最远点对，得4；\textbf{平均链接}取所有跨簇点对距离的平均，得$(4+3)/2=3.5$。同一个“簇间距离”因为定义不同会给出不同数值，后续合并顺序也可能改变。

层次聚类从单个样本开始，不断合并相近的簇，或者从全体开始反向分割。凝聚式层次聚类通过连接准则决定两个簇之间的距离，例如单链接、完全链接和平均链接。它可以用树状图展示不同尺度的群体结构，但并不是无条件优于 K-means。

层次聚类适用于中小数据和希望观察分组层级的场景。对百万量级目录，直接构造完整的样本间距离可能不现实；数据量较大时需要考虑近邻图、抽样或更适合的算法。

\subsection{两群之间的距离怎样定义}
凝聚层次聚类从每条记录各成一簇开始，每次合并距离最近的两簇，直到形成完整合并层次。关键不只是点间距离，还包括\textbf{簇间连接规则}。单连接取两簇之间最近的一对点距离；完全连接取最远距离；平均连接对跨簇点对距离取平均。Ward方法选择使簇内平方误差增加最少的合并，通常与欧氏几何对应。

\worked{同一对群体，不同连接可以得出不同相似程度}
一维群体 $A=\{0,1\}$，$B=\{3,6\}$。四个跨簇距离为3、6、2、5。单连接距离为2，完全连接为6，平均连接为4。单连接容易沿近邻桥形成长链，完全连接更限制整体直径。三者不是对同一标准的不同近似，而是使用不同群体相似定义。

Ward合并增加量可写为
\[
\Delta_{A,B}=\frac{n_An_B}{n_A+n_B}\|\boldsymbol\mu_A-\boldsymbol\mu_B\|^2.
\]
它来自将两组各自均值替换为共同均值后的平方误差分解。例如 $A=\{0,1\}$、$B=\{3,4\}$，两个均值0.5与3.5、各2条，增加量为 $2\times2/4\times3^2=9$。原组内总SSE为1，合并均值2后SSE为10，确实增加9。

\subsection{树状图与决策树不是同一种树}
层次聚类的树状图记录样本群体怎样合并，横向叶端是记录或初始群，合并高度与所用距离或代价有关。第8章决策树的节点则是用于预测的特征阈值，两者不能因为都画成树就混淆。

在某高度切开层次图可以形成平面分组，但高度意义取决于连接规则与尺度。并列距离或轻微扰动可能改变部分合并顺序。普通层次聚类通常也没有像监督分类器那样统一的 \texttt{predict}新样本接口；新对象如何归组需要另行规定，或重新进行聚类。

\section{DBSCAN 与密度连接}

\textbf{DBSCAN}是一种基于局部密度形成群体的方法。先规定半径$\varepsilon$：一个点的$\varepsilon$邻域包含到它距离不超过该半径的点，\textbf{在本书使用的sklearn约定中计入点自身}。邻域内记录数达到\texttt{min\_samples}的称为\textbf{核心点}；不是核心点、但落在某个核心点邻域内的称为\textbf{边界点}；不能归入任何这样的群体则记为\textbf{噪声点}。簇由核心点之间的密度连接扩展，并吸收附近边界点。\textsuperscript{[22]}

用人工一维点$0,0.10,0.20,0.35,2.00$，设$\varepsilon=0.16$、\texttt{min\_samples=3}。点0.10的邻域含0、0.10、0.20，因此是核心点；0.20的邻域含0.10、0.20、0.35，也是核心点。点0和0.35各只有两个邻域记录，却在核心点邻域内，因此是边界点；点2.00只有自身，成为噪声。于是前四点形成一个簇，最后一点标为$-1$。

这一计算说明，DBSCAN不是要求簇内任意两点的距离都小于$\varepsilon$。0和0.35之间距离大于0.16，仍可通过中间核心点形成的结构属于同一簇。

DBSCAN 通过邻域半径 $\varepsilon$ 与最少邻居数 \texttt{min\_samples} 定义密度连通结构。足够密集的区域形成簇，稀疏区域可能被标记为噪声（通常标签为$-1$）。它不需要预先指定簇数，可以发现某些不规则形状，但对不同密度簇、特征尺度和参数仍然敏感。

\fig{ch14_cluster_shapes.pdf}{人工双月与团块数据上的 K-means 和 DBSCAN 对比，用于展示不同算法关于结构形状的假设。}

DBSCAN 输出的噪声标签不等于“测量错误”或“天文学上的特殊天体”。在低密度真实星群、边缘候选和观测精度变化较大的区域，噪声标签可能包含科学上合理的对象。需要通过外部信息与稳定性分析复核。

\subsection{核心、边界和噪声的判定}
设半径 $\epsilon$，点 $i$ 的邻域包含距离不超过 $\epsilon$ 的记录。按本章与sklearn的约定，计数包含点自身。若邻域至少有 \texttt{min\_samples} 条记录，称为\textbf{核心点}。距离某核心点足够近、但自身邻域不够密的点，可以作为\textbf{边界点}；既不是核心点、又不能从核心结构到达的点标为噪声。

\worked{七个一维点的密度分组}
取 $0,0.1,0.2,1.0,1.1,1.2,3.0$，$\epsilon=0.15$，最小样本数3。点0.1的邻域有0、0.1、0.2，因此为核心；两端0与0.2各只有自身和中间点，是边界。同样，1.1为第二个核心，1.0和1.2为其边界。3.0仅有自身，标为噪声。结果形成两个三点簇和一个噪声点。

现在考虑 $0,1,2,3$，半径1.1、最小样本数3。点1和2是核心，彼此连接，端点0和3分别附着。四点形成一个簇，尽管最远两点距离3大于半径1.1。\textbf{半径限制局部连接，不是整个簇的最大直径。}这正使DBSCAN能够沿弯曲结构传播。

\subsection{为什么不能穿过任意稀疏桥继续扩张}
核心点之间可以通过连续核心邻接形成结构；边界点可以被吸纳，但不会像核心一样无限继续扩张。因此一条仅由稀疏边界点构成的桥，不一定把两个密群合并。若一个边界点同时靠近两个核心簇，归属可能随处理顺序而变化；核心结构比个别边界标签更值得稳定性检查。

DBSCAN不必事先指定簇数，但仍需选择距离、半径和密度阈值。相同全局半径面对密度差异很大的真实群体可能顾此失彼：对稀疏群太小会碎裂，对密群外背景太大又会合并。自动得到几个簇不代表模型已经免除参数选择。

\subsection{尺度如何影响密度}
把一个坐标单位放大，邻域球在物理含义上就改变。自行、视差等不同量纲需要明确尺度；简单按样本标准差缩放可能把真实大散布方向压缩，也可能受到背景主导。距离尺度可以来自测量误差或科学上可接受差异，但必须写清依据，并做敏感性比较。

\section{尺度、参数和结果稳定性}

\textbf{轮廓系数}提供一种分组内部诊断。对一个非单点簇中的样本，记$a$为它到同簇其他点的平均距离，$b$为它到其他各簇的平均距离中的最小值，计算$s=(b-a)/\max(a,b)$。当$a=1,b=4$时，$s=0.75$，表示它相对更靠近本簇；若$a=4,b=1$，则$s=-0.75$，提示当前分组与距离结构不协调。它依赖距离和簇划分，不直接使用物理身份；只有一个簇等情形也不满足其计算条件。

对于距离类聚类，特征标准化会改变群体的定义。一个自行变量取值范围大、另一个视差变量范围小，标准化前后可能得到不同聚类。参数也会影响结构：K-means 的 $K$、DBSCAN 的 $\varepsilon$ 和 \texttt{min\_samples} 不应只凭看图选择，还要考虑测量精度、样本密度及研究目的。

结果稳定性可以通过多次初始化、适度扰动输入、重采样和参数范围比较来检查。稳定不是证明正确的充分条件：某个仪器系统误差也可以形成稳定聚类。对专业问题，应同时检查外部物理信息和已知数据质量。

\subsection{轮廓系数衡量的是某种几何紧凑性}
对一个样本，$a$为它与同簇其他样本的平均距离，$b$为它与其他簇的平均距离中最小者，轮廓系数为
\[
s=\frac{b-a}{\max(a,b)}.
\]
若 $a=1,b=4$，则 $s=0.75$；若 $a=4,b=1$，则 $s=-0.75$，说明它与另一个簇平均更近。单点簇的处理采用软件约定，不能凭一个孤立点得到无穷好的分组。

轮廓系数偏好同簇平均距离小、异簇平均距离大的表示。一个长而弯曲的真实群体，沿曲线相隔很远的点会增大同簇距离，所以即使它符合生成身份，轮廓分也可能不高。这不是指标出错，而是指标与目标可能不一致。

\subsection{用外部标签比较时避免编号陷阱}
聚类标签0与参考类别0没有天然对应。可以使用调整兰德指数（ARI）等基于样本对关系的指标，它对簇编号置换不敏感，并对随机一致性作调整。ARI为1表示两种分组一致；接近0表示按其调整基准没有更多一致性，也可能出现负值。外部指标需要参考分组，不能在没有真实身份的科学目录上凭空计算。

稳定性可以通过重复初始化、轻微测量扰动或重采样观察。对重采样，只比较两次共同出现的对象，或在固定参照对象上建立明确映射；不能把长度不同、对象不同的标签数组直接比较。稳定也不自动等于物理真实，一个由强仪器偏差形成的群体可能非常稳定。

\section{生成数据与隐藏标签数据的比较}

\textbf{生成数据}是按明确的数学规则由程序构造的记录。本节的团块数据由几个中心附近随机分散的点组成；双月数据由两条弯曲半环附近的点组成。它们分别由scikit-learn的\texttt{make\_blobs}和\texttt{make\_moons}生成。参数\texttt{noise}控制人为扰动的强度，不是某一台望远镜的测量误差；所有图均标为机制实验。

课堂先使用 \texttt{make\_blobs()} 展示紧凑群体，再用 \texttt{make\_moons()} 展示弯曲结构。把两个算法用在相同数据上，讨论差异如何反映模型假设。随后隐藏 Iris 或 Wine 的标签，在未使用标签的情况下完成聚类，最后再用原有类别作事后对照。

\begin{lstlisting}[caption={两种聚类方法在相同输入上的比较}]
from sklearn.datasets import make_moons
from sklearn.preprocessing import StandardScaler
from sklearn.cluster import KMeans, DBSCAN

X, _ = make_moons(n_samples=350, noise=0.07, random_state=42)
X = StandardScaler().fit_transform(X)
y_k = KMeans(n_clusters=2, n_init=10, random_state=42).fit_predict(X)
y_d = DBSCAN(eps=0.22, min_samples=5).fit_predict(X)
print("K-means groups:", len(set(y_k)))
print("DBSCAN groups including noise:", len(set(y_d)))
\end{lstlisting}

\fig{ch14_kmeans_k.pdf}{团块生成数据上的簇内平方和与簇数的关系。“肘部”只是一种辅助诊断，不自动证明真实簇数。}

\subsection{一个指标与身份不一致的实际例子}
本轮人工生成350条双月记录，噪声标准差参数0.07，种子42。固定K均值为2簇；DBSCAN使用半径0.22、最小样本数5；单连接层次聚类取2簇。生成标签只用于事后核对。结果为：
\begin{center}\small
\begin{tabular}{lrrr}\toprule
方法&簇数&相对生成标签的ARI&轮廓系数\\\midrule
K均值&2&0.2565&0.4888\\
DBSCAN&2&1.0000&0.3264\\
单连接层次聚类&2&1.0000&0.3264\\\bottomrule
\end{tabular}\end{center}
\fig{v3_ch14_K_means.pdf}{人工双月数据的K均值分组。较高轮廓系数反映其紧凑分区偏好，不表示它更接近生成时的两个弯曲群体。}
\fig{v3_ch14_DBSCAN.pdf}{同一人工数据的DBSCAN分组。此设置恢复生成群体，但不能据此推广为所有真实数据的最优方法。}

这组结果非常适合解释“评价必须对应目标”。若任务是找两个紧凑区域，K均值结果可以有意义；若任务是恢复此生成实验的两个月牙身份，DBSCAN和单连接更一致。不能把轮廓系数最高与科学身份最正确当作同一句话。

\subsection{半径扫描显示结构怎样合并}
在相同数据、最小样本数5下，半径0.10得到5簇和17噪声点，0.14得到2簇和4噪声点，0.18得到2簇和1噪声点，0.22与0.30得到2簇且无噪声，0.50把所有记录连成1簇。参数变化可以改变簇数和噪声比例，不只是移动几个边界样本。

真实无标签分析不能用未知身份挑出0.22。应结合距离分布、科学尺度、参数稳定区间和外部信息提出选择，并报告结果对参数的敏感性。在隐藏鸢尾花类别的练习中，同样应先固定分析方案，最后再打开物种列检查；若据物种反复调参，就把它明确称为有参考标签的选择实验。

\section{Gaia 星团邻域实践}

\label{sec:gaia-intro}
\textbf{Gaia}是欧洲空间局的天体测量任务，本案例采用其第三次数据发布（DR3）的目录作为可复现数据版本，不把它写成随时间自动变化的“最新数据”。\textbf{星团}是需要结合位置、运动等证据研究的一组恒星；本练习先寻找某个已知星团周围的候选成员，而不直接证明它们在引力上束缚。\textsuperscript{[15,24]}

目录中每条选定记录对应一个天体。赤经、赤纬描述天空方向，常以度表示；\textbf{视差}$\varpi$描述因观测位置变化而产生的表观角位移，目录常用毫角秒（mas）；\textbf{自行}描述天体在天空方向上的年变化，两个分量通常记为$\mu_{\alpha*}$与$\mu_\delta$，单位为mas/yr。这些角运动量不直接等于以km/s为单位的空间速度，不能把不同单位的列不加检查地相加。\textsuperscript{[15,24]}

可提出的具体问题是：一些候选恒星是否在自行和视差上形成相近群体？颜色--星等图则把一个颜色（两个波段星等之差）和一个亮度量放在两轴上，提供未用于当前聚类的额外一致性检查。质量标记和每项测量的误差也应随数据保留。

Gaia 提供星位、视差、自行和测光等丰富的天体测量量，可用于研究星团邻域的候选成员。课程先准备一个具体数据发布版本的邻域 CSV，记录天区范围、质量筛选和字段单位。在聚类前决定使用哪些特征，是否引入观测误差，是否要对不同物理量进行缩放。

可以先在 $(\mu_{\alpha *},\mu_\delta,\varpi)$ 等运动学/视差特征空间寻找群体。算法输出后，再检查天空分布、颜色--星等图、天体测量误差以及在不同参数下的候选稳定性。不要直接从三维特征聚类得出“物理上束缚的恒星群体”的结论；共同运动学结构可能包含前景或背景污染，也可能受样本选取影响。

\scientificnote{如果教师尚未提供真实 Gaia CSV，则只能完成生成数据与标准数据实验；不得将人为生成的星团样本称为 Gaia 实测结果。专业项目需要保存 Gaia 发布版本、查询条件与可复现的下载步骤。}

\subsection{从邻域目录到可解释的距离}
Gaia星团实践不是“把所有数值列送入DBSCAN”。先定义天区范围和对象质量，再选用于聚类的自行分量、视差等量。自行描述天体在天球上的角位置随时间变化，视差与观测几何有关，二者单位和误差不同。它们的聚集可以支持共同运动候选，但还需要空间分布、颜色—星等关系等证据。

数据中一个标识对应一个源，标识应保持整数或字符串精度，不能为了组成数值矩阵而转成浮点测量特征。某些对象缺少可靠天体测量时，应记录筛选造成的样本限制。用相对误差筛选也会偏向某些距离和亮度范围，分析时不能假设被保留记录完全代表全体。

可以先在开发天区中检查各输入的分布与测量尺度，指定一组变换与参数，再对相近设置比较候选成员集合的重合比例。颜色—星等图若没有参与聚类，可以作为附加一致性检查；若据它不断修改参数，它也已经进入选择过程，应如实记录。

\subsection{科学结论需要比算法标签多一步}
DBSCAN的噪声编号通常为 $-1$，不表示对象没有科学价值；可能是稀有成员、背景或测量较差。某个密度群也不一定就是一个星团，可能来自选择范围、扫描规律或混合总体。候选表应包含源标识、输入、误差、质量、参数设置与外部检查，而不只给出一个簇号。

本书提供真实目录的字段与流程要求，实际数据获取失败时不填入虚构的成员数。离线算法验证使用明确生成的数据，真实目录项目则由学生或教师准备后按附录接口运行。二者分别检验计算机制和科学适用性，不能互相替代。

\chapterpractice{必做：使用 make\_blobs 与 make\_moons 比较 K-means/DBSCAN；隐藏 Wine 标签后执行聚类，再事后对照；专业迁移：在教师提供的 Gaia 邻域 CSV 上做候选聚类和外部信息检查。}
\chapterexercise
\begin{enumerate}
\item 为什么 K-means 的簇内平方和随着 $K$ 增大通常不增加，但不能单独用来选真实簇数？
\item DBSCAN 的噪声标记与真正的异常数据之间有哪些区别？
\item 一个群体在多次初始化下稳定，是否能证明其为真实物理星团？还需要什么证据？
\item 若将自行的单位改变而未调整标准化，距离类聚类的结果可能怎样变化？
\item 为什么用于评估监督模型的 PCA 必须放在训练折内拟合，而只做描述性聚类绘图时可以使用完整待探索数据？
\end{enumerate}

\section*{深化练习与方法检验}
本组题强调解释、计算与验证；D编号在本章内独立计数。计算题应写出所用定义和中间步骤，实验题保留实际输出，包括不支持预期的结果。
\begin{enumerate}[label=D\arabic*.]
\item 对一维点0、1、4、5、9、10及初始中心0、4，手算Lloyd算法到稳定；与中心2.5、9.5对应的分组比较目标。为什么前者稳定而不是最优？
\item 某次聚类的轮廓系数较高，但与已知生成分组的ARI较低，是否逻辑矛盾？使用本章双月数据的结果解释两指标在衡量什么。
\item 在DBSCAN中，一个边界点同时邻接两个不同簇的核心点，但自身不是核心点，它能否像核心点链那样把两簇合并？说明边界归属对输入顺序的可能敏感性。
\item 用Gaia视差和自行找出一簇后，再用同样的视差和自行说它“通过了独立验证”，问题在哪里？提出两项没有直接用于聚类的检查，并说明仍需考虑什么。
\end{enumerate}
